
%
% Group addresses by affiliation; use superscriptaddress for long
% author lists, or if there are many overlapping affiliations.
% For Phys. Rev. appearance, change preprint to twocolumn.
% Choose pra, prb, prc, prd, pre, prl, prstab, prstper, or rmp for journal
%  Add 'draft' option to mark overfull boxes with black boxes
%  Add 'showkeys' option to make keywords appear
\documentclass[aps,prl,twocolumn,groupedaddress,nofootinbib]{revtex4-2}
\usepackage{graphicx}
\usepackage{amsmath}
%\documentclass[aps,prl,preprint,superscriptaddress]{revtex4-2}
%\documentclass[aps,prl,reprint,groupedaddress]{revtex4-2}

% You should use BibTeX and apsrev.bst for references
% Choosing a journal automatically selects the correct APS
% BibTeX style file (bst file), so only uncomment the line
% below if necessary.
%\bibliographystyle{apsrev4-2}


% Helpful cosmology math commands
\newcommand{\Oi}{\Omega_i}
\newcommand{\Oio}{\Omega_{i,0}}
\newcommand{\Om}{\Omega_m}
\newcommand{\Omo}{\Omega_{m,0}}
\newcommand{\Ob}{\Omega_m}
\newcommand{\Obo}{\Omega_{b,0}}
\newcommand{\Or}{\Omega_r}
\newcommand{\Oro}{\Omega_{r,0}}
\newcommand{\Ok}{\Omega_k}
\newcommand{\Oko}{\Omega_{k,0}}
\newcommand{\Ol}{\Omega_\Lambda}
\newcommand{\Olo}{\Omega_{\Lambda,0}}
\newcommand{\Hp}{\mathcal{H}}
\newcommand{\der}[2][]{\frac{d#1}{d#2}}
\newcommand{\dd}[2][]{\frac{d^2#1}{d#2^2}}

\begin{document}

% Use the \preprint command to place your local institutional report
% number in the upper righthand corner of the title page in preprint mode.
% Multiple \preprint commands are allowed.
% Use the 'preprintnumbers' class option to override journal defaults
% to display numbers if necessary
%\preprint{}

%Title of paper
\title{Codes for the $\Lambda$CDM \\ \normalsize AST5220 Final Project}

% repeat the \author .. \affiliation  etc. as needed
% \email, \thanks, \homepage, \altaffiliation all apply to the current
% author. Explanatory text should go in the []'s, actual e-mail
% address or url should go in the {}'s for \email and \homepage.
% Please use the appropriate macro foreach each type of information

% \affiliation command applies to all authors since the last
% \affiliation command. The \affiliation command should follow the
% other information
% \affiliation can be followed by \email, \homepage, \thanks as well.
\author{Simon Silverstein (\#707626)}
% \email[]{simon.silverstein@fys.uio.no}
%\homepage[]{Your web page}
%\thanks{}
%\altaffiliation{}
\affiliation{University of Oslo}

%Collaboration name if desired (requires use of superscriptaddress
%option in \documentclass). \noaffiliation is required (may also be
%used with the \author command).
%\collaboration can be followed by \email, \homepage, \thanks as well.
%\collaboration{}
%\noaffiliation

\date{\today}

\begin{abstract}
    We present code for parameter estimation for $\Lambda$CDM cosmology. \textit{NOTE: This project}\footnote{This project was completed without writing or coding assistance from Large Language Models (LLMs). LLMs were used exclusively for solving issues with PDF compilation.} \textit{and its codes}\footnote{\href{https://github.com/simonmso/cosmosim}{https://github.com/simonmso/cosmosim}} \textit{are a work in progress. Milestones II, III, and IV to come. }
\end{abstract}

% insert suggested keywords - APS authors don't need to do this
%\keywords{}

%\maketitle must follow title, authors, abstract, and keywords
\maketitle

% body of paper here - Use proper section commands
% References should be done using the \cite, \ref, and \label commands
\section{Milestone I}
\subsection{Introduction}
The leading theory of cosmology is the $\Lambda$CDM model, which describes an expanding universe whose main components are (in order of energy density today) dark energy, matter, radiation, and curvature. Constraining the evolution of this universe requires solving a set of differential equations and fitting the solutions to observed data. In this milestone, we present an overview of the theory for the evolution of a homogeneous, isotropic universe. We introduce tests that a code describing this evolution would have to pass, and demonstrate that our code passes them. Finally, we perform an MCMC simulation of the theory for supernova luminosity distances and present the resulting parameter estimates.
\subsection{Theory\footnote{Many of these equations have been lifted directly from the AST5220 course page (\href{https://cmb.wintherscoming.no}{https://cmb.wintherscoming.no}) as recommended.}}
The foundational observation of cosmology is the expansion of the universe, driven by its energy content and composition. We track the expansion with scale factor $a(t)$, defined so $a(t_{\rm today})\equiv a_0\equiv 1$, and observe also the rate of expansion $\dot{a}$ and the \textit{Hubble parameter} $H\equiv\dot{a}/a$. The Friedmann equation governs this expansion; for an energy density $\rho$, the expansion evolves as
\begin{equation}
    \label{eq:fried_basic}
    H^2=\frac{8\pi G}{3}\rho,
\end{equation}
where $G$ is the gravitational constant. $\rho$ is further broken down as the sum of components $\rho=\sum \rho_i$ (for $i$ of matter, radiation, etc.). (\ref{eq:fried_basic}) may be rearranged so that for a given expansion $H$, a \textit{critical density} $\rho_c=3H^2/8\pi G$ of energy is necessary to drive it. It is then simpler to describe the composition of the universe in terms of what fraction of the critical density $\Omega_i=\rho_i/\rho_c$ each component takes up.

The components we consider in this milestone are radiation $\Or$, which includes neutrinos (since the energy of a neutrino is dominated by relativistic momentum), matter $\Om$, which includes all other matter familiar to us as well as dark matter, and dark energy $\Ol$, an energy density that is constant across all $a$. We also treat any overall curvature of the universe as an energy fraction
\begin{equation*}
    \Omega_{k,0} = -\frac{kc^2}{H_0^2}
\end{equation*}
for curvature parameter $k$ since it affects the evolution of the universe similarly. Each of these components evolves with $a$ as
\begin{align}
    \begin{split}
        \label{eq:Omegas}
        \Omega_{k}(a)       & = \frac{\Omega_{k,0}}{a^2H(a)^2/H_0^2}     \\
        \Omega_{m}(a)       & = \frac{\Omega_{m,0}}{a^3H(a)^2/H_0^2}     \\
        \Omega_r(a)         & = \frac{\Omega_{r,0}}{a^4H(a)^2/H_0^2}     \\
        \Omega_{\Lambda}(a) & = \frac{\Omega_{\Lambda 0}}{H(a)^2/H_0^2},
        % \Omega_{\rm CDM}(a) &= \frac{\Omega_{\rm CDM 0}}{a^3H(a)^2/H_0^2} \\
        % \Omega_b(a) &= \frac{\Omega_{b 0}}{a^3H(a)^2/H_0^2} \\
        % \Omega_\gamma(a) &= \frac{\Omega_{\gamma 0}}{a^4H(a)^2/H_0^2} \\
        % \Omega_{\nu}(a) &= \frac{\Omega_{\nu 0}}{a^4H(a)^2/H_0^2} \\
    \end{split}
\end{align}
with subscript 0 indicating the value today. Our Friedmann equation then takes the form
\begin{equation}
    \label{eq:friedmann}
    H=H_0\sqrt{\Omo a^{-3} + \Oro a^{-4} + \Oko a^{-2} + \Olo}.
\end{equation}
We often use instead the unitless parameter $h$, defined so that $H=h\times(100\times\rm km/s/Mpc)$. Since $a$ spans many orders of magnitude from the birth of the universe to today, we do much of our numerical work with the quantity $x \equiv \ln a$. It is also useful when working with experimental data to use the redshift $z(t)$, the fractional change in wavelength of a photon emitted at $t$ and observed today. It is related to $a$ as
\begin{equation}
    1+z = \frac{a_0}{a(t)}.
\end{equation}

Maybe the first question anyone ever asked about the universe is ``how big is it?" This is of course meaningless; the only reasonable question we can answer is `how big is the bit we can see?' Consider a photon emitted at the beginning; the size of the observable universe is the distance that photon travelled since then, known as the \textit{horizon size} or \textit{conformal distance} $\eta$. This is complicated by the fact that earlier on the universe was smaller by factor $a$; for our calculation this is the same as the photon moving $1/a$ times faster. The infinitesimal distance $d\eta$ is then velocity times infinitesimal time $d\eta=(c/a)dt$, or
\begin{equation*}
    \der[\eta]{t} = \frac{c}{a(t)}.
\end{equation*}
Using the definition of $H$, introducing $\Hp\equiv aH$, and changing variables, this is equivalent to
\begin{equation}
    \label{eq:eta}
    \der[\eta]{x} = \frac{c}{\Hp}.
\end{equation}
From the definitions of $H$ and $x$ it is also easy to obtain
\begin{align}
    \label{eq:t}
    \frac{dt}{dx} = \frac{1}{H}.
\end{align}
These two differential equations are valuable. By evolving them numerically (or for special cases solving them analytically), we get $a(t)$, $\eta(t)$, $H(t)$, and derivatives thereof; all the essentials needed to describe the composition of energy over the course of our universe's history.

Often we care about light that was emitted sometime between the beginning and now. For light emitted at time $t$ and observed by us, the \textit{comoving distance} is
\begin{align*}
    % \label{eq:chi}
    \chi(t) = \eta_0 - \eta(t)
\end{align*}
(\textit{comoving} because it is constant between points moving along with the expanding universe). Note that this does not take curvature into account; with curvature, the line-of-sight distance (or \textit{comoving radius}) $r$ between us and the light source is
\begin{align*}
    % \label{eq:r}
    r & = \begin{cases}
              \chi \cdot \frac{\sin \left( \sqrt{|\Omega_{k 0}}| H_0 \chi /c \right)}{\left(\sqrt{|\Omega_{k 0}|}H_0\chi / c\right)}
               & \Omega_{k 0} < 0  \\
              \chi
               & \Omega_{k 0} =0   \\
              \chi \cdot \frac{\sinh \left( \sqrt{|\Omega_{k 0}|} H_0 \chi / c\right)}{\left(\sqrt{|\Omega_{k 0}|}H_0\chi / c\right) }
               & \Omega_{k 0} > 0.
          \end{cases}
\end{align*}
Our models are often fit using `standard candles' which have near-identical luminosity, allowing us to calculate distances from flux measurements. At a \textit{luminosity distance} $d_L$ for a luminosity $L$, the observed flux is $L$ divided by the surface area at that distance $F=L/4\pi d_L^2$. Today, $d_L$ is just the comoving radius $r$; in the past these standard candles seem \textit{further} from us than they are since the expansion of the universe has further decreased the flux, and so we must divide by the scale factor to obtain
\begin{align}
    \label{eq:dL}
    d_L = \frac{r}{a}.
\end{align}
Through the course of this project, we use a fiducial cosmology as a point of comparison. We use the parameter estimates from \textit{Planck 2018} \cite{aghanim_planck_2020}, most importantly\footnote{The value for $\Oro$ is equivalent to CMB temperature $T_{\rm CMB}=2.725$ K and effective massless neutrino number $N_{\rm eff}=3.046$}
\begin{align*}
    \Omo & = 0.313                 \\
    \Oko & = 0                     \\
    \Oro & = 9.346\times 10^{-5}   \\
    \Olo & = 0.6869                \\
    H_0  & = 66.88 {\rm km/s/Mpc}.
\end{align*}


\subsection{Tests}
Although the relationship between $H$ and $a$ (and also $x$) is nontrivial, one can construct scenarios in which $H$ and values derived from it have simple behavior. The easiest of these is when one component of energy dominates the Friedmann equation (\ref{eq:friedmann}), that is $\Oio \to 1$. In the limits of matter, radiation, and dark energy domination, the quantities $\Hp'/\Hp$, $\Hp''/\Hp$, and $\eta H/c$ all have simple behavior (here $\circ'\equiv\der{x}\circ$). We therefore use these quantities to test our code's implementation of $H$, $\Hp$, and their derivatives.

To see what this simple behavior is, consider first the limit of $H$, given by Friedmann equation (\ref{eq:friedmann}). In the limits, (\ref{eq:friedmann}) goes to
\begin{align*}
    \Om \to 1 & : H(a)\to H_0\sqrt{\Omo}a^{-3/2} \\
    \Or \to 1 & : H(a)\to H_0\sqrt{\Oro}a^{-2}   \\
    \Ol \to 1 & : H(a)\to H_0\sqrt{\Olo}.
\end{align*}
Test (A) of our code is the quantity $\frac{1}{\Hp}\der[\Hp]{x}=\frac{\Hp'}{\Hp}$. We differentiate $\Hp(a)=aH(a)$ using the chain rule, divide by $\Hp(a)$, and apply the limits above to obtain
\begin{align}
    \begin{split}
        \label{eq:a}
        \Om \to 1 & : \frac{\Hp'}{\Hp} \to -\frac{1}{2} \\
        \Or \to 1 & : \frac{\Hp'}{\Hp} \to -1           \\
        \Ol \to 1 & : \frac{\Hp'}{\Hp} \to 1.
    \end{split}
\end{align}
Doing so again gives us test (B), for which the quantity $\frac{1}{\Hp}\dd[\Hp]{x}=\frac{\Hp''}{\Hp}$ goes to
\begin{align}
    \begin{split}
        \label{eq:b}
        \Om \to 1 & : \frac{\Hp''}{\Hp} \to 1            \\
        \Or \to 1 & : \frac{\Hp''}{\Hp} \to -\frac{1}{4} \\
        \Ol \to 1 & : \frac{\Hp''}{\Hp} \to 1.
    \end{split}
\end{align}
Our last test is (C), where we note that in the radiation dominated era, the conformal time may be approximated $\eta(x)\approx c/\Hp(x)$. This gives us the nice result that for
\begin{equation}
    \label{eq:c}
    \Or \to 1: \frac{\eta(x) \Hp(x)}{c} \to 1.
\end{equation}

To bookend the periods in which radiation, matter, or dark energy dominate, we mark the times at which two components are equal. For radiation-matter equality, we equate the expressions for $\Or$ and $\Om$ given in (\ref{eq:Omegas}) and solve for $x$, giving us
\begin{equation}
    \label{eq:rm}
    x_{rm}=\ln{\frac{\Oro}{\Omo}}.
\end{equation}
We do the same procedure for matter and dark energy to obtain
\begin{equation}
    \label{eq:mlam}
    x_{m\Lambda}=\frac{1}{3}\ln{\frac{\Omo}{\Olo}}.
\end{equation}
Since the fiducial cosmology has negligible curvature, we do not bother to implement expressions for curvature-component equality. Together, (\ref{eq:rm}) and (\ref{eq:mlam}) are also a useful test of our implementation of $\Oi(x)$.

One similar quantity is the point at which the expansion of the universe begins to accelerate. This happens long after the era of radiation domination, so in calculating it we use the simplification $\Or\approx0$. From there, we differentiate the Friedmann equation with respect to time to obtain
\begin{equation}
    \label{eq:accel}
    x_{\rm accel} = \frac{1}{3}\ln \left(\frac{1}{2}\frac{\Omo}{\Olo}\right).
\end{equation}
\begin{figure}
    \includegraphics[width=\linewidth]{../results/milestone1/domination.png}
    \caption{\label{fig:tests}(Top plot) Share of total radiation $\Or$, matter $\Om$, and dark energy $\Ol$ with respect to $x$ as the fiducial universe evolves, including analytic $x$-values for radiation-matter (RM) and matter-dark energy (M$\Lambda$) equalities (dashed lines). The latter three plots compare the code's output to the analytic limits for tests (A), (B), and (C) in the text (eq.s \ref{eq:a}, \ref{eq:b}, and \ref{eq:c}; marked by horizontal dashed lines).}
\end{figure}

Figure \ref{fig:tests} shows the energy breakdown of the fiducial universe during its evolution, along with the value of our three test quantities. Note first that qualitatively, the analytic radiation-matter and matter-dark energy equalities fall where their respective curves overlap. This is exactly as expected. Further, the three test quantities approach the analytic limits from (\ref{eq:a}), (\ref{eq:b}), and (\ref{eq:c}) respectively. Together, these demonstrate the successful implementation of $\Oi$, $H$, $\Hp$, and their derivatives in the code.

\subsection{Parameter Estimation}
We measure cosmological parameters using a Markov chain Monte Carlo (MCMC) simulation to estimate the probability distribution of a dataset of 31 type Ia supernovae \cite{betoule_improved_2014}. Each supernova has measured redshift $z$ and luminosity distance $d_L$ with uncertainties. We use the \verb|emcee| python package \cite{foreman-mackey_emcee_2013} to calculate the MCMC chain. We assume the errors $\{\sigma_{d_L}\}$ to be Gaussian-distributed, and therefore choose the logarithm of the likelihood to be
\begin{equation}
    \label{eq:log_likely}
    \log p(d_L|\theta)\propto-\frac{1}{2}\frac{(d_{L} - \hat{d}_{L})^2}{\sigma_{d_L}^2}
    + \log(2\pi \sigma_{d_L}^2),
\end{equation}
where $\hat{d}_L$ is the predicted luminosity distance given parameters $\theta$, and the total log-likelihood is the sum of log-likelihoods from each supernova. We use a uniform prior with the logarithm
\begin{equation}
    \label{eq:log_prior}
    \log p(\theta)\propto
    \begin{cases}
                & 0.1 < h_0 < 1.5,    \\
        0       & 0 \leq \Omo \leq 1, \\
                & -1 \leq \Oko \leq 1 \\
        -\infty & \text{otherwise.}
    \end{cases}
\end{equation}
We derive $\Olo$ from $\sum\Oio=1$. Each chain starts at $\theta_{\rm start}=(h_0,\Omo,\Oko)=(0.7,0.3,10^{-5})$, with $\Oro$ fixed at the fiducial $\Oro = 8.532\times 10^{-5}$.

As recommended by the package authors, we estimate the integrated autocorrelation time $\tau$ and only consider chains run for over $50\tau$ to be trustworthy. We also discard the first $\approx\tau$ samples as `burn-in,' so that our final distributions are not biased towards the starting position.

To calculate $\hat{d}_L$, we evolve (\ref{eq:eta}) and (\ref{eq:t}) numerically from very early times ($a\approx10^{-8}$) with the initial conditions $\eta(x_{\rm start})=c/\Hp(x_{\rm start})$ and $t(x_{\rm start})=1/2H(x_{\rm start})$. We then fit a polynomial spline to the result so that we may calculate $\eta(x)$ for any $x$.
\subsection{Results}
\begin{table}
    \caption{\label{tab:results1}Summary of parameter estimates from MCMC on supernova data. $\theta_{\rm best}$ is the value with the highest likelihood, with 68\% of samples falling on the interval $(\theta_{\rm best} - \sigma_-, \theta_{\rm best} + \sigma_+)$.}
    \begin{ruledtabular}
        \begin{tabular}{l r}
            $\theta$                  & $\theta_{\rm best}$ $\sigma_+/-\sigma_-$ \\
            \hline
            \footnotesize{Sampled parameters}                                    \\
            $h_0$                     & 0.7013 +0.0065/-0.0060                   \\
            $\Omega_{m,0}$            & 0.2559 +0.1104/-0.1122                   \\
            $\Omega_{k,0}$            & 0.0800 +0.2806/-0.2752                   \\
            \hline
            \footnotesize{Derived parameters}                                    \\
            $\Omega_{\Lambda,0}$      & 0.6640 +0.1709/-0.1740                   \\
            $\Omega_{m,0}h_0^2$       & 0.1259 +0.0561/-0.0559                   \\
            $\Omega_{\Lambda,0}h_0^2$ & 0.3266 +0.0900/-0.0887                   \\
        \end{tabular}
    \end{ruledtabular}
\end{table}
\begin{table}
    \vspace{1em}
    \begin{ruledtabular}
        \caption{\label{tab:best}Summary of features of the most-likely cosmology.}
        \begin{tabular}{l r}
            Feature                      & Value                       \\
            \hline
            Age $t(x=0)$                 & 13.67 (Gyr.)                \\
            Conformal time $\eta(x=0)/c$ & 47.16 (Gyr.)                \\
            \hline
            \footnotesize{Radiation-Matter equality}                   \\
            $x_{\rm rm}$                 & -8.010                      \\
            $z_{\rm rm}$                 & 3010                        \\
            $t_{\rm rm}$                 & 6.520$\times10^{-5}$ (Gyr.) \\
            \hline
            \footnotesize{Acceleration begins}                         \\
            $x_{\rm accel.}$             & -0.5489                     \\
            $z_{\rm accel.}$             & 0.7313                      \\
            $t_{\rm accel.}$             & 7.203 (Gyr.)                \\
            \hline
            \footnotesize{Matter-Dark Energy equality}                 \\
            $x_{\rm m\Lambda}$           & -0.3178                     \\
            $z_{\rm m\Lambda}$           & 0.3741                      \\
            $t_{\rm m\Lambda}$           & 9.609 (Gyr.)                \\
        \end{tabular}
    \end{ruledtabular}
\end{table}

The results of the MCMC simulation are summarized in Table \ref{tab:results1}. Features of the highest-likelihood cosmology are summarized in Table \ref{tab:best}. Figure \ref{fig:h0} shows the distribution of samples for $h_0$, including the fiducial value of $h_0=0.6688$. The two are in complete disagreement; notice that the fiducial value falls well outside the entire distribution of samples. This is rather alarming. We choose not to speculate on why these two disagree (there are plenty of more qualified people paid to do so), and note only that the fiducial value was measured from a \textit{significantly} earlier point in our universe's history than the supernovae, meaning there is plenty of room for new theories.

Figure \ref{fig:2dhist} shows the distribution for parameters $\Omo$ and $\Olo$. Unlike the distribution for $h_0$, the highest-likelihood values are in close agreement with the fiducial cosmology. However, these parameters are much more poorly constrained, with 95\% intervals that span $\pm100\%$ of the highest-likelihood value. We therefor cannot say that more sensitive tests would not also cause disagreement here.
\begin{figure}
    \includegraphics[width=\linewidth]{../results/milestone1/h0_hist.png}
    \caption{\label{fig:h0}Histogram of samples of $h_0$ from the Monte Carlo fit to supernova luminosity distances ($n_{\rm bins}=50$). The solid red line marks the highest likelihood value, the dotted lines the 18th and 82nd percentiles (values reported in table \ref{tab:results1}). The fiducial value of $h_0=0.6688$ (black line) is included for reference.}
\end{figure}
\begin{figure}
    \includegraphics[width=\linewidth]{../results/milestone1/omega_hist_MLam.png}
    \caption{\label{fig:2dhist}Histogram of samples of $\Olo$ and $\Omo$ from the Monte Carlo supernovae fit ($n_{\rm bins}=100\times100$). The red `x' marks the sample with the highest likelihood, and 68\% and 95\% of the samples fall within the inner and outer dashed lines respectively (values reported in table \ref{tab:results1}). The fiducial value of $\Omo=0.313, \Olo\approx0.69$ (black star) is included for reference. }
\end{figure}

Figure \ref{fig:supernovae} plots $d_L$ for our best cosmology alongside the supernova data to which it was fit. At these low redshifts, there is little qualitative difference between the best cosmology and the fiducial cosmology. The fact that for all universes $d_L$ is linear with $z$ in the low-$z$ limit is the reason for our poor constraining ability.
\begin{figure}
    \includegraphics[width=\linewidth]{../results/milestone1/supernovea.png}
    \caption{\label{fig:supernovae}Luminosity distance $d_L(z)$ vs. redshift $z$ for the fiducial cosmology and for the highest likelihood cosmology from MCMC simulations (`Best'; see table \ref{tab:results1}). Measures of $d_L(z)$ with experimental uncertainty are plotted for the 31 supernovae fit by the MCMC simulation. }
\end{figure}

Figure \ref{fig:eta} reports the conformal time $\eta/c$ for our cosmology alongside a `toy' cosmology of $\Omo=0.5,\Olo\approx0.5,T_{\rm CMB}=2.725$ K for comparison. This same `toy' cosmology is included in Figure \ref{fig:Hp}, which demonstrates that although the expansion of the universe has long been slowing down, it has relatively recently begun accelerating again.
\begin{figure}
    \includegraphics[width=\linewidth]{../results/milestone1/eta.png}
    \caption{\label{fig:eta}Conformal time $\eta(x)/c$ for the best cosmology from the MCMC simulation, the fiducial cosmology, and, for comparison, a `toy' cosmology of $\Omo=0.5,\Olo\approx0.5,T_{\rm CMB}=2.725$ K. Note that the `Best' line lies almost entirely beneath the fiducial.}
\end{figure}

\begin{figure}
    \includegraphics[width=\linewidth]{../results/milestone1/Hp.png}
    \caption{\label{fig:Hp}$\Hp(x)$ for the best cosmology from the MCMC simulation, the fiducial cosmology, and, for comparison, a `toy' cosmology of $\Omo=0.5,\Olo\approx0.5,T_{\rm CMB}=2.725$ K. Note that the `Best' line lies almost entirely beneath the fiducial.}
\end{figure}

Finally, Figure \ref{fig:t} shows our numerical solution to (\ref{eq:t}) for the best, fiducial, and toy cosmologies.

\begin{figure}
    \includegraphics[width=\linewidth]{../results/milestone1/t.png}
    \caption{\label{fig:t}Age of the universe $t(x)$ for the best cosmology from the MCMC simulation, the fiducial cosmology, and, for comparison, a `toy' cosmology of $\Omo=0.5,\Olo\approx0.5,T_{\rm CMB}=2.725$ K. The current age of the universe $t(0)$ for the `best' model is $t=13.67$ Gyr., as reported in table \ref{tab:results1}. Note that the `Best' line lies almost entirely beneath the fiducial.}
\end{figure}


Broadly, these parameter estimates support the prevailing theory of a universe dominated by a cosmological constant $\Lambda$, with the remaining energy density today coming from matter (remember that we fixed $\Oro$ to the value derived from CMB temperature measurements). We cannot rule out curvature, though it is most likely to be negligible. With the limited data we have here, we are unable to distinguish between baryonic and dark matter.


\section{Milestone II}
\subsection{Introduction}
The primary concern of the second milestone is recombination, the period of the universe's evolution during which the hot plasma of early on cooled enough to allow hydrogen atoms to form without being immediately reionized. Without enough energy to ionize hydrogen, the photons of this time have been travelling mostly uninterrupted since, and are what we today observe as the cosmic microwave background (CMB), one of the key experimental observations of modern cosmology and the focus of Milestone IV. Knowing when and how recombination occured is a necessary step towards modelling these observables.

As this project is the work of a Master's student, we neglect helium in our treatment of recombination. Following big bang nucleosynthesis, approximately 22\% of the baryon mass was tied up in $^4$He (tk cite dodelson p 95), a significant enough fraction that all the following results should be taken with a rather over-sized grain of salt. Further, we neglect reionization, which took place around $z\approx7.7$ (tk cite planck) and would appear as a jump in the free electron fraction $X_e$ with a corresponding hump in the visibility function $\tilde g$ at that time.

\subsection{Theory}
On its journey from a source to an observer, light passing through some medium loses intensity to effects like scattering and absorbtion. We quantify this by relating the observed intensity $I$ to the source intensity $I_0$ with $I=I_0e^{-\tau}$, and refer to $\tau$ as the `optical depth.' For the cosmological case of photons scattering off of free electrons in an expanding universe, $\tau$ is defined to be
\begin{equation}
    \tau(\eta) = \int_{\eta}^{\eta_0} n_e \sigma_T a d\eta'.
\end{equation}
While useful, it is simpler computationally to solve the differential form
\begin{equation}
    \frac{d\tau}{dx} = -\frac{c n_e \sigma_T }{H},
    \label{eq:tau}
\end{equation}
where $\sigma_T$ is the Thompson cross-section, $H(x)$ is known from Milestone 1, and $n_e(x)$ is the number density of free electrons at a given time.

Using $\tau$, we can compute the \textit{visibility function} $\tilde g(x)$ with
\begin{equation}
    \tilde{g}(x) = -\tau' e^{-\tau}.
\end{equation}
Since by construction $\int_{-\infty}^0 \tilde g dx = 1$, $\tilde g(x)$ is a probability density, and we interpret it as the probability a photon observed today was last scattered at time $x$. The peak of this function is the time at which it is most likely a photon was last scattered, aptly named the `time of last scattering.'

The difficulty in computing $\tau(x)$ comes from determining the free electron density $n_e$ as a function of $x$. Early on, when temperatures are high, the dominant interaction $\gamma + H \leftrightarrow e^-+p$ is in equilibrium, and the Boltzman equation for the interaction leads to the relation
\begin{equation}
    \frac{n_en_p}{n_H} = \frac{n_e^{(0)}n_p^{(0)}}{n_H^{(0)}},
    \label{eq:Saha_frac}
\end{equation}
Here $n_i^{(0)}$ is the number density of species $i$ at zero chemical potential (see tk dodelson tk callin). To make this useful, we introduce the \textit{fractional electron density} $X_e$ as\footnote{Note that this follows tk Callin and is slightly different from the definition which tk Dodelson and others use for the \textit{free electron fraction} $X_e = n_e / (n_e + n_H)$. Although under $n_H=n_b$ ours is functionally the same, upon including He (which we here exclude) ours would no longer represent `the fraction of electrons which are free,' allowing for percentages greater than 1 and causing confusion to AST5220 students everywhere.}
\begin{equation}
    X_e= \frac{n_e}{n_H} = \frac{n_e}{n_b},
    \label{eq:Xe}
\end{equation}
which under (\ref{eq:Saha_frac}) gives
\begin{equation}
    \label{eq:Saha}
    \frac{X_e^2}{1-X_e} = \frac{1}{n_b} \left(\frac{m_e
        T_b}{2\pi}\right)^{3/2} e^{-\epsilon_0/T_b},
\end{equation}
known as the \textit{Saha equation} (as is (\ref{eq:Saha_frac})) [tk cite saha]. With
\begin{equation}
    n_H = n_b \approx \frac{\Omega_{b0} \rho_{c0}}{m_H a^3},
    \label{eq:nH approx}
\end{equation}
this equation can be solved analytically by the quadratic formula.

It is only for $X_e > 0.99$ that the equilibrium approximation (\ref{eq:Saha_frac}) holds; afterwards we switch to solving the \textit{Peebles} ODE (tk cite peebles)
\begin{equation}
    \frac{dX_e}{dx} = \frac{C_r(T_b)}{H} \left[\beta(T_b)(1-X_e) - n_H
        \alpha^{(2)}(T_b)X_e^2\right],
    \label{eq:peebles}
\end{equation}
which takes into account ionization from the ground state (through $\beta(T_b)$), recombination to excited states (through $\alpha^{(2)}(T_b)$), and the chance that an excited state atom may be reionized before decaying to the ground state (through $C_r(T_b)$). See tk cite Ma for details on these factors.

Together, the Saha and Peebles equations allow us to solve for $n_e$, from which we compute $\tau(x)$ and $\tilde g(x)$.

A final quantity of interest for this milestone is the sound horizon $s(x)$, which describes the farthest a sound wave could have traveled through the photon-baryon plasma, and therefor the length scale of correlations in the CMB at recombination and the matter power spectrum today. It is given in differential form by
\begin{equation}
    \frac{ds(x)}{dx} = \frac{c_s}{\mathcal{H}}
    \label{eq:sound_horizon}
\end{equation}
with $c_s = c \sqrt{\frac{R}{3(1+R)}}$, $R = 4\Omega_{\gamma 0}/3\Omega_{b 0} a$, and $s(x_{\rm ini}) = c_s(x_{\rm ini})/\mathcal{H}(x_{\rm ini})$.
\subsection{Results}
All results for Milestone II are generated from our `toy' cosmology of $h_0=0.7$, $\Obo=0.05$, $\Omega_{\rm{CDM},0}=0.45$, $\Olo\approx0.5$, and $T_{\rm{CMB}}=2.7255$ K.

\begin{figure}[h]
    \includegraphics{../results/milestone2/X_e_z.png}
    \caption{\label{fig:Xe}Fractional electron density $X_e$ computed once using only the Saha approximation (eq. \ref{eq:Saha}; dashed line) and again with the Peebles ODE (\ref{eq:peebles}) for $X_e<0.99$ (solid line).}
\end{figure}
Figure \ref{fig:Xe} reports the fractional electron density $X_e$. The `Saha + Peebles' solution shows the Saha approximation (\ref{eq:Saha}) until $X_e < 0.99$, after which the numerical solution to the Peebles ODE (\ref{eq:peebles}) is used with the Saha approximation as its initial condition. As evidenced by the figure, the Saha approximation diverges from the Peebles ODE around $z\approx1250$. For this cosmology, the freezeout abundance today is $X_e(x=0)=1.919\times10^{-4}$, as reported in Table \ref{tab:results2}.

\begin{figure}[h]
    \includegraphics{../results/milestone2/tau.png}
    \caption{\label{fig:tau}Numerical solution in the `toy' cosmology to the optical depth $\tau$ with first and second derivatives w.r.t. $x=\ln(a)$.}
\end{figure}
\begin{figure}[h]
    \includegraphics{../results/milestone2/g.png}
    \caption{\label{fig:g}Visibility function $\tilde g$ in the `toy' cosmology with first and second $x$-derivatives (separated into three plots for legibility). By definition, the scale of $\tilde g$ is such that $\int_{-\infty}^0 \tilde g dx = 1$.}
\end{figure}
Using our solution to $X_e$ and relation (\ref{eq:nH approx}), we numerically integrated the $\tau$ differential equation (\ref{eq:tau}). Figure \ref{fig:tau} shows the results of this. Notice the sharp drop in opacity around $x\approx-7$; this is when recombination takes place, and is in line with the exponential decrease in $X_e$. Figure \ref{fig:g} shows the visibility function $\tilde g$ with a corresponding peak at the same time, implying that a background photon observed today is almost guaranteed to have last scattered during recombination (again, this is not including reionization). Exact times for recombination and last scattering are reported in Table \ref{tab:results2}.

As a test of our implementation of $\tilde g$, a quick numerical integration differed from the expected value of 1 by no more than $10^{-5}$.

Finally, after numerically integrating (\ref{eq:sound_horizon}), Table \ref{tab:results2} reports the sound horizon at recombination to be 125.01 Mpc. We should therefor expect this length scale to be imprinted on the CMB power spectrum in the coming milestones.

With these solutions, we now have the tools to compute the perturbations of Milestone III.
\begin{table}[h]
    \caption{\label{tab:results2}Important quantities from Milestone II}
    \begin{ruledtabular}
        \begin{tabular}{l r r}
            \footnotesize{Time of last scattering (peak of $\tilde g$)}                         \\
            $x_{\rm{ls}}$                                       &        & -7.160               \\
            $z_{\rm{ls}}$                                       &        & 1285.5               \\
            $t_{\rm{ls}}$ (yr.)                                 &        & 251453               \\
            \hline
            \footnotesize{Time of recombination ($X_e = 0.1$):} & Full   & Saha                 \\
            $x_{\rm{rc}}$                                       & -7.141 & -7.142               \\
            $z_{\rm{rc}}$                                       & 1262.0 & 1263.4               \\
            $t_{\rm{rc}}$ (yr.)                                 & 258999 & 258550               \\
            \hline
            \footnotesize{Freezeout abundance}                                                  \\
            $X_e(x=0)$                                          &        & $1.919\times10^{-4}$ \\
            \hline
            \footnotesize{Sound horizon at recombination}                                       \\
            $s(x_{\rm{rc}})$ (Mpc)                              &        & 125.01               \\
        \end{tabular}
    \end{ruledtabular}
\end{table}
\section{Milestone III}
\subsection{Introduction}
While sitting and reading this report, it may occur to you that the universe today is not the isotropic and homogeneous fluid we have so far treated it as. We are surrounded by stuff, sheep and busses and planets and galaxies and parts of the cosmic microwave background which are fractionally cooler or hotter than the average. Milestone III concerns itself with the development of this stuff (especially the latter), with the goal of allowing us in Milestone IV to predict the statistical properties of the CMB anisotropies and matter power spectrum.

Our modeling is based on first-order perturbations to the quantities with which we are already familiar, such as the density $\rho_i$ of each component, the temperature $T$, and the metric $g$. The introduction of perturbations changes the Boltzman equation for each component, which together with the Einstein equations leaves us a set of coupled differential equations to solve numerically. These are reproduced in tk \ref{appendix}.
\subsection{Theory}
As the derivation of the system of differential perturbation equations is rather involved, we merely introduce its characters and give a brief qualitative sketch of how one might solve for them. We begin with the metric potentials $\Psi$ and $\Phi$ in the \textit{Conformal Newton gauge}
\begin{align*}
    g_{00} = & -1-2\Psi                \\
    g_{ij} = & a^2\delta_{ij}[1+2\Phi]
\end{align*}
where the delta function ensures that the off-diagonal elements are $g_{0i}=g_{i0}=0$, as with the unperturbed metric, and the time- and space-like components still have opposite sign. $\Phi$ and $\Psi$ together encode perturbations to the scale factor $a$ and to spacetime itself.

Next are density perturbations for cold dark matter $\delta_{\rm{CDM}}$ and  for baryons $\delta_b$. These are defined fractionally, so that
\begin{equation*}
    \delta_i\equiv\frac{\delta\rho_i}{\rho_i}
\end{equation*}
or, in terms of the number density,
\begin{equation*}
    n_i=\bar{n}_i[1+\delta_i]
\end{equation*}
where $\bar n_i$ is the average number density for species $i$. Similarly, we define the photon temperature fluctuation $\Theta$ so that
\begin{equation*}
    T=\bar T[1+\Theta].
\end{equation*}
We work with $\Theta$ primarily through its multipole expansion cite tk
\begin{equation*}
    \Theta_\ell\equiv\frac{i^\ell}{2}\int_{-1}^{1}d\mu\mathcal{P}_\ell(\mu)\Theta(\mu),
\end{equation*}
where $\mathcal{P}_\ell(\mu)$ is the order-$\ell$ Legendre polynomial.

The basic process of solving for each perturbation starts by expanding a component's distribution function $f(\mathbf{x}, \mathbf{p}, t)$ to first order in the perturbation. This is then plugged in to the appropriate Boltzman equation, including a collision term for the dominant interactions for the component, resulting in an expression involving both  time and spacial derivatives of the perturbation. By transforming the entire expression to Fourier space, spacial derivatives become factors of wavevector $k_i$, and our partial differential equations become coupled ordinary differential equations. Better yet, each Fourier mode is independent, and may be individually solved for each desired $k$.

Expanding $f$ and the collision terms to first order in the perturbations leaves one more perturbation to keep track of: the group velocity $v_i$ for a fluid. Specifically, we consider the velocity of cold dark matter $v_{\rm{CDM}}$ and of baryons $v_b$.

There are two interactions of note. Coulomb scattering $e^- + p \leftrightarrow e^- + p$ couples electrons and protons (our only baryons since we aren't including He and heavier elements) at all times\footnote{It is for this reason we do not evolve seperate $e^-$ and $p$ perturbations}. Compton scattering $e^-+\gamma \leftrightarrow e^- +\gamma$ couples the photons to free electrons (and thus to baryons) before recombination. These have further qualitative effects, such as the baryon velocity $v_b$ influencing the photon dipole $\Theta_1$.

Dark matter does not interact with $\delta_b$ or $\Theta$ by definition. However, all three components do influence the metric perturbations through gravity as governed by the Einstein equations
\begin{equation*}
    G^\mu_\nu=8\pi GT^\mu_\nu.
    \label{eq:einstein}
\end{equation*}
Metric perturbations show up in the Einstein tensor $G^\mu_\nu$, which ultimately depends only on the metric $g$. Baryons, cold dark matter, radiation, and the cosmological constant (not yet mentioned in this milestone since we assume it to be constant, i.e. unperturbed) all contribute to the energy-momentum tensor $T^\mu_\nu$. In return, $\Phi$ and/or $\Psi$ show up in the expressions for the evolution of every other component.

One other prediction to come out of equation \ref{eq:einstein} is that the sum $\Phi+\Psi$ is proportional to
\begin{equation}
    \Phi+\Psi \propto \frac{a^2}{k^2}\delta_\gamma\Theta_2
    \label{eq:pot_cancel}
\end{equation}
(neglecting neutrinos and polarization) and therefore $\Phi=-\Psi$ for most times and scales (tk cite dodelson p 146).

With our set of differential equations, the final necessity for successful solving is a set of initial conditions. These are given by inflationary physics, and described in tk ref dodelson. Our specific initial conditions are version of these modified in tk cite callin to give values a nonzero starting point for numerical convergence.

\subsection{Numerical evolution}
There is one small hitch in numerically evolving the coupled set of differential equations (Appendix tk \ref{appendix}) for the perturbations. A term proportional to
\begin{equation}
    \tau'[3\Theta_1 + v_b]
    \label{eq:unstable}
\end{equation}
appears in the expressions for both $\Theta_1'$ and $v_b'$. From Milestone II we know that $|\tau'|$ grows exponentially large at early times. On the other hand, $\Theta_1$ and $v_b$ are nearly zero initially, and (\ref{eq:unstable}) is therefor unstable at early times. We refer to the period in which this is an issue as the `tight coupling regime,' and remedy it by solving modified expressions for $\Theta_1'$ and $v_b'$ (also included in appendix tk \ref{appendix}) which are numerically stable.

In theory, the proper point at which to switch between integrating the tight equations to the full equations depends on the wavenumber $k$ for which you are integrating (larger tktk ). In practice, however, this point is essentially `close to but before recombination' for all but the largest $k$. For this reason, we always switch at $x=-8.3$, regardless of $k$.

\subsection{Results}
We numerically evolved the perturbation equations for 300 values of $k$ ranging logarithmically from $k\eta_0=1.0$ to $k\eta_0=3000.0$ tk check this and from $x=-tk$ to $x=0$. From these, we demonstrate the behavior for perturbations at $k=0.2$/Mpc, $k=0.02$/Mpc, and $k=0.002$/Mpc. All results in this milestone are generated for the `toy' cosmology from Milestone II.
\begin{figure}[h]
    \includegraphics[width=\linewidth]{../results/milestone3/density.png}
    \caption{\label{fig:density}Numerical solution for density perturbations for the `toy' cosmology, including for radiation $\delta_\gamma=4\Theta_0$. $\delta_i$ are plotted for three different magnitudes of $k$.}
\end{figure}
Figure \ref{fig:density} displays the results for density perturbations. As expected, all perturbations in the early universe are suppressed while they lie outside the horizon. The smaller wavelengths (larger $k$) enter earlier, with the specific point at which they cross the horizon dictating their behavior. The first mode to enter (dotted line) does so before recombination, and as such, the baryon perturbation $\delta_b$ (tk change label to be abs) closely follows the photon perturbation $\delta_\gamma$. This is the high rate of Compton scattering at play; only at recombination ($x\approx-7.1$) do the two split, with baryons afterwards tracing $\delta_{\rm{CDM}}$.

Comparatively, the largest scale (small $k$) perturbation doesn't cross the horizon until long after recombination, by which point $\delta_b$ and $\delta_\gamma$ are completely independent, leaving $\delta_b$ to perfectly trace $\delta_{\rm{CDM}}$.

In the error of matter-domination, $\delta_{\rm{CDM}}$ and $\delta_b$ grow exponentially, as matter overdensities attract more matter, further growing the overdensity.

\begin{figure}[h]
    \includegraphics[width=\linewidth]{../results/milestone3/velocity.png}
    \caption{\label{fig:velocity}Velocity perturbations for cold dark matter and baryons (upper) and for radiation (lower, separated for legibility) for three different $k$, where $v_\gamma = -3\Theta_1$. For clarity, $v_\gamma$ are desaturated after oscillations look dense ($x = tk$ for $k=0.02 /$Mpc and $x = $ for $k=0.2 /$Mpc).}
\end{figure}

Figure \ref{fig:velocity} shows similar photon-baryon behavior for the velocity perturbations, with $v_b$ for all three wavenumbers closely following $v_\gamma$ up until recombination. Since we define $v_\gamma$ as proportional to the dipole $\Theta_1$, we interpret this behavior as the photon dipole inducing a bulk baryon velocity, with electrons traveling against the photon temperature gradient facing more energetic interactions than those traveling with it.

Notice that unlike matter, the radiation perturbations $\delta_\gamma = 4\Theta_0$ and $v_\gamma =    -3\Theta_1$ do not grow larger in magnitude in the matter-dominated era. This is to be expected, since in this period the expansion of the universe far outpaces Compton scattering. Also, CDM and baryons experience gravitational forces to a much more significant degree, causing perturbations to grow where they don't for photons.

In contrast to the density perturbations, which are constant in the radiation-dominated era, the velocity perturbations grow exponentially from the beginning. tk why is this.

\begin{figure}[h]
    \includegraphics[width=\linewidth]{../results/milestone3/theta2.png}
    \caption{\label{fig:theta2}Photon quadrupole evolution for three different wavenumbers $k$. For clarity, both larger $k$ are greyed after oscillations become dense ($x = tk$ for $k=0.02 /$Mpc and $x = $ for $k=0.2 /$Mpc).}
\end{figure}

Figure \ref{fig:theta2} demeonstrates the typical behavior of photon moments beyond the first (the monopole $\Theta_0$) and second (the dipole $\Theta_1$). These higher moments are suppressed by the large value of $\tau'$ in the tight coupling regime, and only begin to oscillate around the period of recombination. Here again we see larger $k$ experiencing more rapid oscillations, though still not growing in baseline or overall amplitude.

\begin{figure}[h]
    \includegraphics[width=\linewidth]{../results/milestone3/phi.png}
    \caption{\label{fig:potentials}Metric perturbations as the universe evolves for three different values of $k$.}
\end{figure}

Finally, Figure \ref{fig:potentials} tracks the metric perturbations for our three different regimes. For large $k$ with small wavelength, $\Phi(k)$ is able to enter the horizon before recombination, and as such decreases substantially. This is the effect of the radiation dominated era's rapid expansion working against the ability of gravity to pull matter together, diminished further by the high pressures and temperatures of the time keeping overdensities from growing too much.

$\Phi(k)$ outside the horizon at this time are unaffected. These enter into the scene during matter-domination, at which point all potentials become nearly constant at whatever amplitude they had before. It is during this period of constant potential that density perturbations grow exponentially. Only at times very close to today, when dark energy dominates, do we see potentials again dropping off.

Note that $|\Phi(k)+\Psi(k)|$ is almost exactly 0 for all but a short period after recombination for small $k$. This is predicted from (\ref{eq:pot_cancel}); the short period in which the quantity is nonzero corresponds exactly to the time after recombination when both the quadrupole $\Theta_2$ has started to grow (Figure \ref{fig:theta2}) and the photon density $\rho_\gamma$ has not yet decreased too much (first plot of Figure \ref{fig:tests}).
\section{Milestone IV}
\subsection{Introduction}
Having solved for the evolution of perturbations in our universe, we now have the information we need to predict the power spectra of matter and of the CMB. This is the concern of Milestone IV, and is broadly speaking as simple as integrating our multipole results $\Theta_\ell(k)$ over $k$. We compare our CMB result directly to the canonical CMB data from Planck 2018 tk cite, and find poor agreement, likely a result of us neglecting neutrinos, polarization, He, and reionization.
\subsection{Theory}
From Milestone III, we have a solution for $\Theta_\ell(k, x)$, yet we are only ever able to measure the temperature at the present ($x=0$)\footnote{For this entire section, we continue to use $x=\ln a$ as a time variable. We remind the reader that this is a different notation from the derivations in tk ref Dodelson, which uses $x$ as a position variable and keeps track of time with $\eta$ or $t$.}. Further, it is typical for experiments to map temperature onto a spherical map, so that temperature $T(\hat n)$ is a function of direction $\hat n$ on the sky. Finally, given that we are limited to measuring $T(\hat n)$ from only one position in space, a proper analysis should compare only the statistics of $T(\hat n)$ to those predicted by $\Theta_\ell(k, x=0)$. These are the motivations for our decomposition of $T$ into spherical harmonics and our use of the correlations $C_\ell$.

We expand $T(\hat n)$ into spherical harmonics like
\begin{equation}
    T(\hat{n}) = \sum_{\ell=1}^\infty \sum_{m=-\ell}^\ell a_{\ell m} Y_{\ell m}(\hat{n})
\end{equation}
for coefficients $a_{\ell m}$ and where $Y_{\ell m}(\hat{n})$ are the eigenfunctions of Lapace's equation on a sphere (and are thus orthogonal to one another). This $\ell$ turns out to be the same as that for $\Theta_\ell$. Since for a given $\ell$, all $a_{\ell m}$ are drawn randomly from the same Gaussian distribution with mean $\langle a_{\ell m} \rangle = 0$, we compare our results to their variance
\begin{equation}
    C_\ell \equiv \langle |a_{\ell m}|^2 \rangle = \langle a_{\ell m}a_{\ell m}^* \rangle.
\end{equation}
In terms of $\Theta_\ell$ then, this `angular power spectrum' is given by
\begin{equation}
    C_\ell = \frac{2}{\pi}\int k^2P_{\rm primordial}(k) \Theta_\ell^2(k)dk,
    \label{eq:c_ell_naive}
\end{equation}
where $P_{\rm primordial}(k)$ is the power spectrum following inflation [tk cite dodelson]. Computationally, (\ref{eq:c_ell_naive}) leaves us with a problem: from Milestone III, we know each expression for $\Theta_\ell$ to be dependent on $\Theta_{\ell - 1}$, and so solving up to $\ell\approx3000$ requires numerically evolving a very large number of coupled differential equations.

The technique of `line integration' solves this for us. Using the fact that we only care about $\Theta_\ell$ at $x = 0$, we can express the $\ell$-temperature moment as an integral over time of some source function $\tilde{S}(k,x)$ times a Bessel function $j_\ell$
\begin{equation}
    \Theta_\ell(k, x=0) = \int_{-\infty}^{0} \tilde{S}(k,x) j_\ell[k(\eta_0-\eta)] dx.
    \label{eq:lineofsight}
\end{equation}
In practice, we integrate (\ref{eq:lineofsight}) from some early time when the integrand is small. The source function is derived from the expressions for the derivatives of each multipole (tk appendix tk), and the terms of $\tilde{S}(k,x)$ may be separated by their physical origin. The dominant contribution is the \textit{Sachs-Wolfe} (SW) term, which is just the monopole $\Theta_0$ weighted by the visibility function $\tilde{g}$, with corrections from $\Psi$ and $\Theta_2$ for energy lost to gravity and a small effect from the angular dependence of Compton scattering
\begin{equation}
    \tilde{S}^{\rm{SW}}(k,x) = \tilde{g}\left[ \Theta_0 + \Psi + \frac{1}{4}\Theta_2\right].
    \label{eq:sw}
\end{equation}
Next we have the \textit{integrated Sachs-Wolfe} term, which acts as correction to the previous term accounting for the time dependence of $\Phi$ and $\Psi$, particularly after recombination when there is still radiation energy density and during the late dark energy-dominated era.
\begin{equation}
    \tilde{S}^{\rm{ISW}}(k,x)=e^{-\tau} \left[\Psi^\prime-\Phi^\prime\right].
    \label{eq:isw}
\end{equation}
A `doppler' term we present without comment, and finally include a `polarization' term, which had we included polarization would have accounted for its effect on the temperature power spectrum
\begin{equation}
    \tilde{S}^{\rm{Doppler}}(k,x)=-\frac{1}{ck}\frac{d}{dx}(\mathcal{H}\tilde{g}v_b)
    \label{eq:doppler}
\end{equation}
\begin{equation}
    \tilde{S}^{\rm{Pol.}}(k,x) = \frac{3}{4c^2k^2} \frac{d}{dx}\left[\mathcal{H}\frac{d}{dx} (\mathcal{H}\tilde{g}\Theta_2)\right].
    \label{eq:pol_s}
\end{equation}
With these four terms, the full source function is
\begin{equation}
    \tilde{S}(k,x) = \tilde{S}^{\rm{SW}} + \tilde{S}^{\rm{ISW}} + \tilde{S}^{\rm{Doppler}} + \tilde{S}^{\rm{Pol.}}.
    \label{eq:source}
\end{equation}
The final element needed for computing $C_\ell$ is then the primordial power spectrum. From inflationary physics, we borrow the Harrison-Zel'dovich (tk cite?) spectrum
\begin{equation}
    \frac{k^3}{2\pi^2} P_{\rm{primordial}}(k) = A_s \left(\frac{k}{k_{\rm{pivot}}}\right)^{n_s-1},
\end{equation}
where we use spectral index $n_s=0.96$ and for a pivot scale of $k_{\rm{pivot}}=0.05$/Mpc have an amplitude $A_s\approx 2\times10^{-9}$. With that, we are ready to calculate $C_\ell$.

We keep the matter power spectrum in terms of wavenumber $k$. We then already have the elements we need to calculate the full spectrum. Since our normalization separates the primordial spectrum from the Fourier components $\Delta_M(k, x)$ we computed, the full spectrum is
\begin{equation}
    P(k,x) = |\Delta_M(k,x)|^2P_{\rm primordial}(k)
    \label{eq:matterpwr}
\end{equation}
where the Fourier components are
\begin{equation}
    \Delta_M(k,x) \equiv \frac{2}{3}\left(\frac{ck}{H_0}\right)^2 \frac{\Phi(k,x)}{\Omega_{M0}a^{-1}}.
\end{equation}
From Milestone III we know the behavior of a mode is largely determined by whether it entered the horizon before or after recombination. Therefor, a very useful quantity to know is the wavenumber $k_{\rm{eq}}$ for perturbations which entered the horizon right at matter-radiation equality. This is
\begin{equation}
    k_{\rm{eq}} = \frac{a_{\rm eq} H(a_{\rm eq})}{c}.
    \label{eq:k_eq}
\end{equation}
We should see considerable difference in the matter power spectrum to the left and right of $k_{\rm{eq}}$.
\subsection{Results}
Aside from the two direct comparisons with data (figures \ref{fig:planck_low} and \ref{fig:planck_high}), all results for this milestone are computed with parameters of the `toy' cosmology from Milestone II.

\begin{figure}[h]
    \includegraphics[width=\linewidth]{../results/milestone4/transfer.png}
    \caption{\label{fig:transfer}Transfer function for the `toy' cosmology scaled by $\sqrt{\ell(\ell - 1)}$ for tk six different values of $\ell$.}
\end{figure}
Figure \ref{fig:transfer} shows the behavior of the transfer function $\Theta_\ell(k)$ following our line of sight integration (\ref{eq:lineofsight}). Each $\Theta_\ell$ has similarly rapid oscillatory behavior which scales approximately with $(\ell(\ell + 1))^{-1/2}$, beginning at higher $k$ and continuing longer for higher $\ell$. These effects are qualitatively inherited from the Bessel functions $j_\ell$. In the simplest approximations, $\tilde{S}(k, x)$ is dominated by the Sachs-Wolfe term, itself approximately a dirac delta function $\tilde{g}(x)\approx\delta(x - x_{\rm{eq}})$ scaled by $\Theta_0$. This delta picks $j_\ell(k(\eta_0-\eta))$ out of integral (\ref{eq:lineofsight}), leading to $\Theta_\ell(k)\propto j_\ell(k(\eta_0-\eta))$ and the rapid oscillations in the figure.

Squaring this and scaling by $1/k$ gives a predictable result for the $C_\ell$ integrand (excluding the primordial spectrum), as demonstrated in Figure \ref{fig:integrand}. We integrate this using simple trapezoid integration, sampled at a frequency of $\Delta k = 2\pi/12\eta_0$ so that there are $\approx 12$ points per oscillation.
\begin{figure}[h]
    \includegraphics[width=\linewidth]{../results/milestone4/integrand.png}
    \caption{\label{fig:integrand}Integrand of the temperature correlation function, scaled by $\ell(\ell + 1)$ and not including the primordial spectrum. Results for six different values of $\ell$ are shown.}
\end{figure}

\begin{figure}[h]
    \includegraphics[width=\linewidth]{../results/milestone4/components.png}
    \caption{\label{fig:components}Contributions to $C_\ell$ for the `toy' cosmology. Each line represents (\ref{eq:c_ell_naive}) integrated with only one of the four source terms, aside from `Full $C_\ell$' (black line) which includes them all.}
\end{figure}
Figure \ref{fig:components} shows the contribution of each of the four source terms to the final power spectrum. This was achieved by replacing $\tilde{S}(k,x)$ in (\ref{eq:c_ell_naive}) with one of $\tilde{S}^{\rm{SW}}$, $\tilde{S}^{\rm{ISW}}$, $\tilde{S}^{\rm{Doppler}}$, or $\tilde{S}^{\rm{Pol.}}$ (equations \ref{eq:sw}, \ref{eq:isw}, \ref{eq:doppler}, or \ref{eq:pol_s}), and integrating as usual.

One by one then: the Sachs-Wolfe term is the largest contribution at almost all scales. It is nearly constant in the low-$\ell$ ($\ell < 30$) regime, since these larger-scale perturbations have only recently entered in to the horizon. We then see a large first peak around $\ell=180$, with expected oscillatory behavior thereafter. Only at small scales does the SW spectrum taper off as diffusion damping kicks in.

The integrated Sachs-Wolf term contributes at very large scales and around the first Sachs-Wolfe peak. This is as expected from the form of (\ref{eq:isw}), which tells us that $\tilde{S}^{\rm{ISW}}$ will only be significant when $\tau\lesssim1$ and the potentials are changing. For the largest scales, $\tau$ is tiny and recent dark-energy domination is causing $\Psi$ and $\Phi$ to diminish, as we learned from figures \ref{fig:tau} and \ref{fig:potentials}. As for the contribution to the first peak, the only other time ISW could contribute is earlier than matter domination (which froze $\Psi$ and $\Phi$) and later than $x\approx-8$ (since otherwise $\tau$ would be too large). This is to say that the only other major ISW contribution is to scales which entered in to the horizon at about recombination, which is precisely what the first peak corresponds to.

The Doppler term is insignificant at large scales, which correspond to perturbations that have only recently entered the horizon and have therefor not Doppler shifted much. It then contributes mostly uniformly to the spectrum until being diffusion damped away.

Regretfully, we do not have a clear or intuitive explanation for the behavior of the `Polarization' term (which, again, does not actually contain any polarization perturbation information in our implementation).

\begin{figure}[h]
    \includegraphics[width=\linewidth]{../results/milestone4/planck_low.png}
    \caption{\label{fig:planck_low}$C_\ell$ for the fiducial cosmology compared to Planck low-$\ell$ data (tk cite planck).}
\end{figure}
Figure \ref{fig:planck_low} compares our computational prediction for $C_\ell$ given the fiducial cosmology to the low-$\ell$ experimental data from the Planck telescope (tk cite Planck). Given the high experimental uncertainty in the region (owing both to the Planck systematics and the high cosmic variance at large scales) it is difficult to say much about the quality of our prediction; we merely remark that by eye, our fit passes through $\approx2/3$ of the experimental uncertainty bars and therefor is experimentally plausible (as any good scientist knows to check).

\begin{figure}[h]
    \includegraphics[width=\linewidth]{../results/milestone4/planck_high.png}
    \caption{\label{fig:planck_high}$C_\ell$ for the fiducial cosmology compared to Planck high-$\ell$ data (tk cite planck). `He fudging' refers to transforming $\ell$ and $C_\ell$ by the fixed factors described in (\ref{eq:fudge})}
\end{figure}
The high-$\ell$ Planck data of Figure \ref{fig:planck_high} is another story entirely. We fail terribly at reproducing the experimental spectrum. This is almost certainly due to our exclusion of reionization, He, polarization, and neutrinos, all of which shift and/or suppress/exagerate certain features. We may `fudge' the shifting from neutrinos and damping from helium with
\begin{align}
    \ell\to   & \ell^{1.018}                       \\
    C_\ell\to & C_\ell\exp(-0.05(\ell/200)^{1.5}),
    \label{eq:fudge}
\end{align}
which as the figure shows much more closely matches the Planck data.

\begin{figure}[h]
    \includegraphics[width=\linewidth]{../results/milestone4/matter.png}
    \caption{\label{fig:matter}Matter power spectrum $P(k,x=0)$ for the present day. $k_{\rm{eq}}$ marks the scale of perturbation that crosses the horizon during recombination.}
\end{figure}
Finally, Figure \ref{fig:matter} shows the matter power spectrum $P(k, x=0)$ for the present time. It is marked by $k_{\rm{eq}}$ calculated from (\ref{eq:k_eq}), with different behavior to each side as predicted. To the left of $k_{\rm{eq}}$, $P(k)$ grows with $k$ since as modes get smaller they enter into the horizon earlier and therefor have more time to collect. This changes at $k_{\rm{eq}}$, marking the point beyond which modes are small enough to have entered into the horizon before recombination. From Figure \ref{fig:potentials}, we know that potentials are suppressed in the radiation-dominated epoch, with larger $k$ affected more. We see this as the drop-off in the matter power spectrum to the right of $k_{\rm{eq}}$.
\section{Conclusion}
We have here demonstrated the moderate success of our code in calculating the angular temperature power spectrum $C_\ell^{TT}$ and matter power spectrum $P(k)$ for a set of cosmological parameters
\clearpage
\appendix
\section{Appendix}
\label{appendix}
Initial Conditions
\begin{align*}
    \Psi             & = -\frac{1}{\frac{3}{2} + \frac{2f_\nu}{5}}                               \\
    \Phi             & = -(1+\frac{2f_\nu}{5})\Psi                                               \\
    \delta_{\rm CDM} & = \delta_b = -\frac{3}{2} \Psi                                            \\
    v_{\rm CDM}      & = v_b = -\frac{ck}{2\mathcal{H}} \Psi                                     \\
                     & \text{Photons:}                                                           \\
    \Theta_0         & = -\frac{1}{2} \Psi                                                       \\
    \Theta_1         & = +\frac{ck}{6\mathcal{H}}\Psi                                            \\
    \Theta_2         & = -\frac{20ck}{45\mathcal{H}\tau^\prime} \Theta_1                         \\
    \Theta_\ell      & = -\frac{\ell}{2\ell+1} \frac{ck}{\mathcal{H}\tau^\prime} \Theta_{\ell-1} \\
\end{align*}

Tight coupling
\begin{align*}
    q=
     & \frac{-[(1-R)\tau^\prime + (1+R)\tau^{\prime\prime}](3\Theta_1+v_b) -\frac{ck}{\mathcal{H}}\Psi}{(1+R)\tau^\prime + \frac{\mathcal{H}^\prime}{\mathcal{H}} -1}                                        \\
     & + \frac{(1-\frac{\mathcal{H}^\prime}{\mathcal{H}})\frac{ck}{\mathcal{H}}(-\Theta_0 +2\Theta_2) - \frac{ck}{\mathcal{H}}\Theta_0^\prime}{(1+R)\tau^\prime + \frac{\mathcal{H}^\prime}{\mathcal{H}} -1} \\
    v_b^\prime=
     & \frac{1}{1+R}                                                                                                                                                                                         \\
     & \times\left[-v_b - \frac{ck}{\mathcal{H}}\Psi + R(q +\frac{ck}{\mathcal{H}}(-\Theta_0 + 2\Theta_2) - \frac{ck}{\mathcal{H}}\Psi)\right]                                                               \\
    \Theta^\prime_1=
     & \frac{1}{3} (q - v_b^\prime)
\end{align*}

Full
\begin{align*}
    \Theta^\prime_0       =
     & -\frac{ck}{\mathcal{H}} \Theta_1 - \Phi^\prime,                                                                      \\
    \Theta^\prime_1      =
     & \frac{ck}{3\mathcal{H}}\left[\Theta_0 - 2\Theta_2 + \Psi\right] + \tau^\prime\left[\Theta_1 + \frac{1}{3}v_b\right], \\
     & 2 \leq \ell < \ell_{\textrm{max}}:                                                                                   \\
    \Theta^\prime_\ell   =
     & \frac{ck}{(2\ell+1)\mathcal{H}}\left[\ell\Theta_{\ell-1} - (\ell+1)\Theta_{\ell+1}\right]                            \\
     & + \tau^\prime\left[\Theta_\ell - \frac{1}{10}\Pi
    \delta_{\ell,2}\right],                                                                                                 \\
     & \ell = \ell_{\textrm{max}}:                                                                                          \\
    \Theta_{\ell}^\prime =
     & \frac{ck}{\mathcal{H}}\left[
        \Theta_{\ell-1}-\frac{\ell+1}{k\eta(x)}\Theta_\ell\right]+\tau^\prime\Theta_\ell,
    \\
\end{align*}

\begin{align*}
    \delta_{\rm CDM}^\prime & = \frac{ck}{\mathcal{H}} v_{\rm CDM} - 3\Phi^\prime                  \\
    v_{\rm CDM}^\prime      & = -v_{\rm CDM} -\frac{ck}{\mathcal{H}} \Psi                          \\
    \delta_b^\prime         & = \frac{ck}{\mathcal{H}}v_b -3\Phi^\prime                            \\
    v_b^\prime              & = -v_b - \frac{ck}{\mathcal{H}}\Psi + \tau^\prime R(3\Theta_1 + v_b) \\
\end{align*}

\begin{align*}
    \Phi^\prime =  & \Psi - \frac{1}{3}\left(\frac{ck}{\mathcal{H}}\right)^2 \Phi                                                                                                                                                         \\
                   & + \frac{H_0^2}{2\mathcal{H}^2}                                                \left[\frac{\Omega_{\rm CDM 0}\delta_{\rm CDM}}{a}  + \frac{\Omega_{b 0}\delta_b}{a} + \frac{4\Omega_{\gamma 0}\Theta_0}{a^{2}}\right] \\
    \Psi         = & -\Phi - \frac{12H_0^2}{c^2k^2a^2}\left[\Omega_{\gamma 0}\Theta_2 + \Omega_{\nu 0}\mathcal{N}_2\right]                                                                                                                \\
\end{align*}


\bibliographystyle{apsrev4-2}
\bibliography{AST5220}


% Put \label in argument of \section for cross-referencing
%\section{\label{}}
% \subsection{}
% \subsubsection{}

% If in two-column mode, this environment will change to single-column
% format so that long equations can be displayed. Use
% sparingly.
%\begin{widetext}
% put long equation here
%\end{widetext}

% figures should be put into the text as floats.
% Use the graphics or graphicx packages (distributed with LaTeX2e)
% and the \includegraphics macro defined in those packages.
% See the LaTeX Graphics Companion by Michel Goosens, Sebastian Rahtz,
% and Frank Mittelbach for instance.
%
% Here is an example of the general form of a figure:
% Fill in the caption in the braces of the \caption{} command. Put the label
% that you will use with \ref{} command in the braces of the \label{} command.
% Use the figure* environment if the figure should span across the
% entire page. There is no need to do explicit centering.

% \begin{figure}
% \includegraphics{}%
% \caption{\label{}}
% \end{figure}

% Surround figure environment with turnpage environment for landscape
% figure
% \begin{turnpage}
% \begin{figure}
% \includegraphics{}%
% \caption{\label{}}
% \end{figure}
% \end{turnpage}

% tables should appear as floats within the text
%
% Here is an example of the general form of a table:
% Fill in the caption in the braces of the \caption{} command. Put the label
% that you will use with \ref{} command in the braces of the \label{} command.
% Insert the column specifiers (l, r, c, d, etc.) in the empty braces of the
% \begin{tabular}{} command.
% The ruledtabular enviroment adds doubled rules to table and sets a
% reasonable default table settings.
% Use the table* environment to get a full-width table in two-column
% Add \usepackage{longtable} and the longtable (or longtable*}
% environment for nicely formatted long tables. Or use the the 
% placement option to break a long table (with less control than 
% in longtable).
% \begin{table}% add  placement to break table across pages
% \caption{\label{}}
% \begin{ruledtabular}
% \begin{tabular}{}
% Lines of table here ending with \\
% \end{tabular}
% \end{ruledtabular}
% \end{table}

% Surround table environment with turnpage environment for landscape
% table
% \begin{turnpage}
% \begin{table}
% \caption{\label{}}
% \begin{ruledtabular}
% \begin{tabular}{}
% \end{tabular}
% \end{ruledtabular}
% \end{table}
% \end{turnpage}

% Specify following sections are appendices. Use \appendix* if there
% only one appendix.
%\appendix
%\section{}

% If you have acknowledgments, this puts in the proper section head.
%\begin{acknowledgments}
% put your acknowledgments here.
%\end{acknowledgments}

% Create the reference section using BibTeX:
% \bibliography{basename of .bib file}

\end{document}
%
% ****** End of file apstemplate.tex ******

